% ===========================================================================
%  13 — Diagramas de atividade
% ===========================================================================
\chapter{Diagramas de atividade}
\label{cap:atividades}

Este capítulo reúne o diagrama de atividade de \textbf{todos} os processos do
sistema --- os automáticos, os manuais e os de operação. A notação é UML de
atividades~\cite{omg2017uml}: círculo cheio para início, círculo com anel para
término, retângulo arredondado para ação, losango para decisão e barra
horizontal para bifurcação ou junção paralela.

\begin{table}[H]
\caption{Índice dos diagramas de atividade}
\label{tab:atividades}
\begin{tabularx}{\linewidth}{@{}L{14mm}XL{22mm}@{}}
\toprule
\headrow \thd{ID} & \thd{Processo} & \thd{Disparo}\\
\midrule
\req{AT-01} & Carga do histórico corporativo (ETL) & Manual\\
\zebra \req{AT-02} & Canonicalização de um rótulo bruto & Interno\\
\req{AT-03} & Treino do motor de diagnóstico & Manual\\
\zebra \req{AT-04} & Indexação de um documento orientativo & Manual ou
  automático\\
\req{AT-05} & Diagnóstico de um novo evento & API\\
\zebra \req{AT-06} & \textit{Guardrail} e geração da prescrição & Interno\\
\req{AT-07} & Conversa com o agente & API\\
\zebra \req{AT-08} & Ingestão de telemetria industrial & Mensagem MQTT\\
\req{AT-09} & Cadastro de documento pela interface & Usuário\\
\zebra \req{AT-10} & Indexação automática no primeiro \textit{boot} &
  \textit{Startup}\\
\req{AT-11} & Consulta a um painel analítico & Usuário\\
\zebra \req{AT-12} & Integração contínua & \textit{Push} ou \textit{pull
  request}\\
\req{AT-13} & Provisionamento e implantação em nuvem & Manual\\
\zebra \req{AT-14} & Ciclo de tratamento de lacuna documental & Negócio\\
\bottomrule
\end{tabularx}
\end{table}

% ---------------------------------------------------------------------------
\section{AT-01 --- Carga do histórico corporativo}

\begin{figure}[H]
\centering
\fitwidth{%
\begin{tikzpicture}[x=1cm, y=1cm]
  \node[actstart] (i) at (3.2,0) {};
  \node[action] (a1) at (3.2,-1.4) {Ler o CSV com tipagem de datas};
  \node[action] (a2) at (3.2,-3.1) {Normalizar os rótulos brutos
    {\scriptsize (minúsculas, sem espaços)}};
  \node[action] (a3) at (3.2,-4.9) {Canonicalizar \textbf{todos} os rótulos
    distintos};
  \node[decision] (d1) at (3.2,-6.9) {Todos casam com alguma regra?};
  \node[actionerr] (e1) at (10.6,-6.9) {Levantar exceção nomeando o rótulo};
  \node[actend] (fe) at (10.6,-8.9) {};
  \node[action] (a4) at (3.2,-9.0) {Deduplicar por \cd{id} e descartar linhas
    incompletas};
  \node[action] (a5) at (3.2,-10.8) {Criar as tabelas, se ausentes};
  \node[action] (a6) at (3.2,-12.5) {Popular \cd{fault\_families} e
    \cd{label\_map}};
  \node[action] (a7) at (3.2,-14.2) {Carregar os eventos de forma idempotente};
  \node[actionok] (a8) at (3.2,-16.0) {Imprimir o resumo: rótulos, famílias e
    linhas carregadas};
  \node[actend] (f) at (3.2,-17.8) {};

  \draw[flow] (i) -- (a1);
  \draw[flow] (a1) -- (a2);
  \draw[flow] (a2) -- (a3);
  \draw[flow] (a3) -- (d1);
  \draw[flow] (d1) -- (e1) node[flowlbl, midway, above=1pt]{não};
  \draw[flow] (e1) -- (fe);
  \draw[flow] (d1) -- (a4) node[flowlbl, pos=0.45, right=2pt]{sim};
  \draw[flow] (a4) -- (a5);
  \draw[flow] (a5) -- (a6);
  \draw[flow] (a6) -- (a7);
  \draw[flow] (a7) -- (a8);
  \draw[flow] (a8) -- (f);

  \node[actnote, text width=48mm, anchor=north west] at (7.6,-10.4)
    {\textbf{Ordem deliberada:} a auditoria completa dos rótulos precede
     qualquer escrita. Se um rótulo novo aparecer, o banco permanece intacto e o
     operador corrige a regra antes de reexecutar.};
  \draw[actnoteline] (7.5,-11.6) -- (5.1,-11.6);
\end{tikzpicture}}
\caption[AT-01 Carga do histórico]{\req{AT-01} --- carga do histórico
  corporativo. O ramo de exceção não é tratamento de erro incidental: é a porta
  de qualidade de dados do sistema (\req{RNF-14}).}
\label{fig:at01}
\end{figure}

% ---------------------------------------------------------------------------
\section{AT-02 --- Canonicalização de um rótulo bruto}

\begin{figure}[H]
\centering
\fitwidth{%
\begin{tikzpicture}[x=1cm, y=1cm]
  \node[actstart] (i) at (3.2,0) {};
  \node[action] (a1) at (3.2,-1.4) {Normalizar: minúsculas e remoção de espaços
    nas bordas};
  \node[action] (a2) at (3.2,-3.2) {Tomar a próxima regra da lista
    \textbf{ordenada}};
  \node[decision] (d1) at (3.2,-5.2) {A expressão casa?};
  \node[decision] (d2) at (3.2,-7.6) {Há mais regras?};
  \node[actionok] (ok) at (10.8,-5.2) {Devolver a família da regra
    {\scriptsize (nome, tipo, descrição)}};
  \node[actionerr] (er) at (10.8,-7.6) {Levantar exceção: rótulo exige curadoria
    humana};
  \node[actend] (f1) at (10.8,-9.4) {};

  \draw[flow] (i) -- (a1);
  \draw[flow] (a1) -- (a2);
  \draw[flow] (a2) -- (d1);
  \draw[flow] (d1) -- (ok) node[flowlbl, midway, above=1pt]{sim};
  \draw[flow] (d1) -- (d2) node[flowlbl, pos=0.5, right=2pt]{não};
  \draw[flow] (d2) -- (er) node[flowlbl, midway, above=1pt]{não};
  \draw[flow] (d2.west) -- (0.6,-7.6) -- (0.6,-3.2) -- (a2.west)
    node[flowlbl, pos=0.5, left=1pt]{sim};
  \draw[flow] (ok) -- (f1);
\end{tikzpicture}}
\caption[AT-02 Canonicalização]{\req{AT-02} --- canonicalização de um rótulo
  bruto. A avaliação é sequencial e a \emph{primeira} regra que casa vence, o
  que torna a ordem das regras parte da especificação: regras de falha precedem
  regras de estado operacional.}
\label{fig:at02}
\end{figure}

% ---------------------------------------------------------------------------
\section{AT-03 --- Treino do motor de diagnóstico}

\begin{figure}[H]
\centering
\fitwidth{%
\begin{tikzpicture}[x=1cm, y=1cm]
  \node[actstart] (i) at (3.4,0) {};
  \node[action] (a1) at (3.4,-1.4) {Carregar e canonicalizar o conjunto de
    dados};
  \node[action] (a2) at (3.4,-3.2) {Construir a matriz: 18 brutas + 5
    derivadas};
  \node[bar, minimum width=40mm] (fk) at (9.0,-4.6) {};
  \node[action, text width=32mm] (v1) at (3.0,-6.4) {\textbf{Avaliação 1}
    \newline \textit{holdout} por sessão};
  \node[action, text width=32mm] (v2) at (9.0,-6.4) {\textbf{Avaliação 2}
    \newline divisão aleatória estratificada};
  \node[action, text width=32mm] (v3) at (15.0,-6.4) {\textbf{Avaliação 3}
    \newline sessão, sem temperatura};
  \node[bar, minimum width=40mm] (jn) at (9.0,-8.4) {};
  \node[action, text width=40mm] (a3) at (9.0,-10.0) {Retreinar padronizador e
    classificador com \SI{100}{\percent} dos dados};
  \node[action, text width=40mm] (a4) at (9.0,-11.9) {Indexar o histórico
    completo na busca por vizinhos};
  \node[actionok, text width=40mm] (a5) at (9.0,-13.8) {Gravar 4 artefatos e
    \arq{metrics.json} com os 3 protocolos};
  \node[actend] (f) at (9.0,-15.6) {};

  \draw[flow] (i) -- (a1);
  \draw[flow] (a1) -- (a2);
  \draw[flow] (a2.south) -- (3.4,-4.0) -- (7.0,-4.0) -- (7.0,-4.53);
  \draw[flow] (7.6,-4.68) -- (3.0,-4.68) -- (v1.north);
  \draw[flow] (9.0,-4.68) -- (v2.north);
  \draw[flow] (10.4,-4.68) -- (15.0,-4.68) -- (v3.north);
  \draw[flow] (v1.south) -- (3.0,-8.32) -- (7.6,-8.32);
  \draw[flow] (v2.south) -- (9.0,-8.32);
  \draw[flow] (v3.south) -- (15.0,-8.32) -- (10.4,-8.32);
  \draw[flow] (jn.south) -- (a3);
  \draw[flow] (a3) -- (a4);
  \draw[flow] (a4) -- (a5);
  \draw[flow] (a5) -- (f);

  \node[actnote, text width=44mm, anchor=north west] at (16.6,-9.6)
    {As três avaliações são independentes e podem ser lidas em paralelo. O
     modelo \emph{servido} não é nenhum dos três: é retreinado depois, com todos
     os dados.};
  \draw[actnoteline] (16.5,-10.6) -- (11.1,-10.6);
\end{tikzpicture}}
\caption[AT-03 Treino]{\req{AT-03} --- treino do motor de diagnóstico. A
  bifurcação paralela expressa a independência entre os três protocolos de
  avaliação; a junção garante que todos concluam antes do treino final.}
\label{fig:at03}
\end{figure}

% ---------------------------------------------------------------------------
\section{AT-04 --- Indexação de um documento orientativo}

\begin{figure}[H]
\centering
\fitwidth{%
\begin{tikzpicture}[x=1cm, y=1cm]
  \node[actstart] (i) at (3.4,0) {};
  \node[action] (a1) at (3.4,-1.4) {Remover \textit{chunks} anteriores do mesmo
    arquivo};
  \node[action] (a2) at (3.4,-3.2) {Extrair texto por página, rastreando a seção
    corrente};
  \node[decision] (d1) at (3.4,-5.3) {Algum \textit{chunk} foi produzido?};
  \node[actionerr] (ocr) at (11.0,-5.3) {Renderizar páginas a
    \SI{150}{dpi} e transcrever por OCR multimodal};
  \node[action] (a3) at (3.4,-7.6) {Acumular \textit{chunks} de
    $\approx$\num{1200} caracteres com sobreposição 150};
  \node[action] (a4) at (3.4,-9.5) {Gerar \textit{embeddings}
    (\cd{RETRIEVAL\_DOCUMENT}, dim.\ 768)};
  \node[decision] (d2) at (3.4,-11.6) {Contagem de vetores = de textos?};
  \node[actionerr] (er) at (11.0,-11.6) {Levantar exceção: desalinhamento
    silencioso};
  \node[action] (a5) at (3.4,-13.9) {Indexar um conjunto por família coberta,
    com metadados};
  \node[actionok] (a6) at (3.4,-15.8) {Registrar documento e cobertura no banco
    {\scriptsize (\textit{best-effort})}};
  \node[actend] (f) at (3.4,-17.6) {};
  \node[actend] (fe) at (11.0,-13.4) {};

  \draw[flow] (i) -- (a1);
  \draw[flow] (a1) -- (a2);
  \draw[flow] (a2) -- (d1);
  \draw[flow] (d1) -- (ocr) node[flowlbl, midway, above=1pt]{não (escaneado)};
  \draw[flow] (ocr.south) -- (11.0,-7.6) -- (a3.east)
    node[flowlbl, pos=0.72, above=1pt]{blocos transcritos};
  \draw[flow] (d1) -- (a3) node[flowlbl, pos=0.5, right=2pt]{sim};
  \draw[flow] (a3) -- (a4);
  \draw[flow] (a4) -- (d2);
  \draw[flow] (d2) -- (er) node[flowlbl, midway, above=1pt]{não};
  \draw[flow] (er) -- (fe);
  \draw[flow] (d2) -- (a5) node[flowlbl, pos=0.5, right=2pt]{sim};
  \draw[flow] (a5) -- (a6);
  \draw[flow] (a6) -- (f);
\end{tikzpicture}}
\caption[AT-04 Indexação documental]{\req{AT-04} --- indexação de um documento.
  O OCR é acionado \emph{apenas} quando a extração direta não produz conteúdo. A
  verificação de contagem de vetores previne a falha mais perigosa da etapa:
  associar o vetor de um trecho ao texto de outro.}
\label{fig:at04}
\end{figure}

% ---------------------------------------------------------------------------
\section{AT-05 --- Diagnóstico de um novo evento}

\begin{landscape}
\begin{figure}[H]
\centering
\fitpage{%
\begin{tikzpicture}[x=1cm, y=1cm]

  % --- raias
  \foreach \y/\h/\nome in {0/2.6/{Interface HTTP}, -2.6/4.6/{Motor de diagnóstico},
                           -7.2/2.6/{Prescrição}, -9.8/2.0/{Persistência}} {
    \draw[lane] (3.2,\y) rectangle (23.0,\y-\h);
    \draw[lanehdr] (0.4,\y) rectangle (3.2,\y-\h);
    \node[font=\scriptsize\bfseries\color{white}, text width=24mm,
          rotate=0] at (1.8,\y-\h/2) {\nome};
  }

  \node[actstart] (i) at (4.4,-1.3) {};
  \node[action, text width=30mm] (a1) at (7.2,-1.3)
    {Validar \textit{payload} contra \cd{SensorEvent}};
  \node[decision, text width=22mm] (d0) at (11.6,-1.3) {Esquema válido?};
  \node[actionerr, text width=28mm] (e0) at (16.2,-1.3)
    {Responder \cd{422} com os campos inconsistentes};
  \node[actend] (f0) at (20.4,-1.3) {};

  \node[action, text width=30mm] (a2) at (7.2,-3.9)
    {Calcular as 5 derivadas e padronizar};
  \node[action, text width=28mm] (a3) at (12.0,-3.9)
    {Classificar: família e probabilidade};
  \node[action, text width=28mm] (a4) at (16.6,-3.9)
    {Recuperar 25 vizinhos e calcular concordância};
  \node[action, text width=28mm] (a5) at (16.6,-6.0)
    {Agregar estatísticas históricas da família};
  \node[action, text width=28mm] (a6) at (12.0,-6.0)
    {Classificar a confiança (produto dos dois sinais)};
  \node[decision, text width=22mm] (d1) at (7.2,-6.0) {É falha?};

  \node[action, text width=30mm] (a7) at (7.2,-8.5)
    {Consultar o \textit{guardrail} de cobertura};
  \node[decision, text width=22mm] (d2) at (11.6,-8.5) {Documentada?};
  \node[actionok, text width=28mm] (a8) at (16.2,-8.5)
    {Recuperar, gerar e citar (AT-06)};
  \node[actionerr, text width=28mm] (a9) at (20.6,-8.5)
    {Aviso + sugestão de registro};

  \node[action, text width=34mm] (a10) at (8.0,-10.8)
    {Persistir auditoria em \cd{diagnoses} \textit{(best-effort)}};
  \node[actionok, text width=30mm] (a11) at (14.0,-10.8)
    {Responder \cd{DiagnoseResponse}};
  \node[actend] (f) at (18.6,-10.8) {};

  \draw[flow] (i) -- (a1);
  \draw[flow] (a1) -- (d0);
  \draw[flow] (d0) -- (e0) node[flowlbl, midway, above=1pt]{não};
  \draw[flow] (e0) -- (f0);
  \draw[flow] (d0.south) -- (11.6,-2.15) -- (7.2,-2.15) -- (a2.north)
    node[flowlbl, pos=0.30, above=1pt]{sim};
  \draw[flow] (a2) -- (a3);
  \draw[flow] (a3) -- (a4);
  \draw[flow] (a4) -- (a5);
  \draw[flow] (a5) -- (a6);
  \draw[flow] (a6) -- (d1);
  \draw[flow] (d1) -- (a7) node[flowlbl, pos=0.5, right=2pt]{sim};
  \draw[flow] (d1.west) -- (4.4,-6.0) -- (4.4,-10.8) -- (a10.west)
    node[flowlbl, pos=0.55, left=1pt]{não (estado)};
  \draw[flow] (a7) -- (d2);
  \draw[flow] (d2) -- (a8) node[flowlbl, midway, above=1pt]{sim};
  \draw[flow] (d2.north) -- (11.6,-7.42) -- (20.6,-7.42) -- (a9.north)
    node[flowlbl, pos=0.30, above=1pt]{não};
  \draw[flow] (a8.south) -- (16.2,-10.15) -- (14.0,-10.15) -- (a11.north);
  \draw[flow] (a9.south) -- (20.6,-10.15) -- (16.3,-10.15);
  \draw[flow] (a10) -- (a11);
  \draw[flow] (a11) -- (f);

\end{tikzpicture}}
\caption[AT-05 Diagnóstico de novo evento]{\req{AT-05} --- diagnóstico de um
  novo evento, em raias por responsabilidade. Note que a persistência da
  auditoria é \textit{best-effort} e paralela ao caminho de resposta: uma falha
  de banco não impede o técnico de receber o diagnóstico.}
\label{fig:at05}
\end{figure}
\end{landscape}

% ---------------------------------------------------------------------------
\section{AT-06 --- \textit{Guardrail} e geração da prescrição}

\begin{figure}[H]
\centering
\fitwidth{%
\begin{tikzpicture}[x=1cm, y=1cm]
  \node[actstart] (i) at (3.4,0) {};
  \node[action] (a1) at (3.4,-1.4) {Obter o nome de exibição da família};
  \node[decision] (d1) at (3.4,-3.4) {Existe \textit{chunk} com esta família?};
  \node[actionerr] (n1) at (11.0,-3.4) {Formatar aviso de falha não documentada
    + sugestão de registro};
  \node[actend] (fe) at (11.0,-5.4) {};
  \node[action] (a2) at (3.4,-5.8) {Montar a consulta: família + sintomas do
    evento};
  \node[action] (a3) at (3.4,-7.6) {Gerar \textit{embedding}
    (\cd{RETRIEVAL\_QUERY})};
  \node[action] (a4) at (3.4,-9.4) {Recuperar \textit{top-k}=6 filtrando pela
    família};
  \node[action] (a5) at (3.4,-11.2) {Montar \textit{prompt} de contexto fechado
    com trechos numerados};
  \node[action] (a6) at (3.4,-13.1) {Gerar com \cd{temperature=0.2} e instrução
    de sistema restritiva};
  \node[actionok] (a7) at (3.4,-15.1) {Extrair citações dos \textbf{metadados}
    dos trechos recuperados};
  \node[actend] (f) at (3.4,-17.0) {};

  \draw[flow] (i) -- (a1);
  \draw[flow] (a1) -- (d1);
  \draw[flow] (d1) -- (n1) node[flowlbl, midway, above=1pt]{não};
  \draw[flow] (n1) -- (fe);
  \draw[flow] (d1) -- (a2) node[flowlbl, pos=0.5, right=2pt]{sim};
  \draw[flow] (a2) -- (a3);
  \draw[flow] (a3) -- (a4);
  \draw[flow] (a4) -- (a5);
  \draw[flow] (a5) -- (a6);
  \draw[flow] (a6) -- (a7);
  \draw[flow] (a7) -- (f);

  \node[actnote, text width=50mm, anchor=north west] at (7.8,-13.4)
    {\textbf{A ordem é a garantia.} A decisão de cobertura ocorre no passo 2, e o
     único caminho até a geração passa pelo ramo \enquote{sim}. Não existe
     percurso de código que chame o modelo para família descoberta.};
  \draw[actnoteline] (7.7,-14.8) -- (5.3,-14.8);
\end{tikzpicture}}
\caption[AT-06 Guardrail e prescrição]{\req{AT-06} --- \textit{guardrail} de
  cobertura e geração da prescrição. As citações são extraídas no penúltimo
  passo, a partir dos metadados do \textit{retrieval} --- nunca do texto
  gerado.}
\label{fig:at06}
\end{figure}

% ---------------------------------------------------------------------------
\section{AT-07 --- Conversa com o agente}

\begin{figure}[H]
\centering
\fitwidth{%
\begin{tikzpicture}[x=1cm, y=1cm]
  \node[actstart] (i) at (3.4,0) {};
  \node[decision] (d0) at (3.4,-1.8) {Agente disponível?};
  \node[action] (fb) at (11.2,-1.8) {Degradar para RAG de um passo com o
    histórico enviado};
  \node[action] (a1) at (3.4,-4.2) {Reiniciar o coletor de citações do
    contexto};
  \node[action] (a2) at (3.4,-6.0) {Recuperar a memória da sessão (esquema
    \cd{adk})};
  \node[action] (a3) at (3.4,-7.8) {Modelo decide a próxima ferramenta};
  \node[decision] (d1) at (3.4,-9.9) {Qual ferramenta?};
  \node[actionsys, text width=32mm] (t1) at (11.4,-7.4)
    {\cd{consultar\_historico}\newline{\scriptsize SQL agregado}};
  \node[actionerr, text width=32mm] (t2) at (11.4,-9.3)
    {\cd{buscar\_documentos}\newline{\scriptsize aplica o \textit{guardrail}}};
  \node[actionsys, text width=32mm] (t3) at (11.4,-11.2)
    {\cd{diagnosticar\_evento}\newline{\scriptsize motor de diagnóstico}};
  \node[actionsys, text width=32mm] (t4) at (11.4,-13.1)
    {\cd{cobertura\_documental}\newline{\scriptsize índice + banco}};
  \node[decision] (d2) at (3.4,-12.6) {Precisa de outra ferramenta?};
  \node[action] (a4) at (3.4,-15.0) {Gerar a resposta final em contexto
    fechado};
  \node[actionok] (a5) at (3.4,-16.9) {Devolver texto + citações do coletor +
    selo do modo};
  \node[actend] (f) at (3.4,-18.8) {};

  \draw[flow] (i) -- (d0);
  \draw[flow] (d0) -- (fb) node[flowlbl, midway, above=1pt]{não};
  \draw[flow] (fb.south) -- (11.2,-16.9) -- (a5.east)
    node[flowlbl, pos=0.82, above=1pt]{resposta e citações};
  \draw[flow] (d0) -- (a1) node[flowlbl, pos=0.5, right=2pt]{sim};
  \draw[flow] (a1) -- (a2);
  \draw[flow] (a2) -- (a3);
  \draw[flow] (a3) -- (d1);
  \draw[flow] (d1.east) -- (9.0,-9.9) -- (9.0,-7.4) -- (t1.west);
  \draw[flow] (9.0,-9.3) -- (t2.west);
  \draw[flow] (9.0,-9.9) -- (9.0,-11.2) -- (t3.west);
  \draw[flow] (9.0,-11.2) -- (9.0,-13.1) -- (t4.west);
  \draw[flow] (t2.south west) -- (7.0,-10.6) -- (7.0,-12.6) -- (d2.east);
  \draw[flow] (d2.west) -- (0.8,-12.6) -- (0.8,-7.8) -- (a3.west)
    node[flowlbl, pos=0.5, left=1pt]{sim};
  \draw[flow] (d2) -- (a4) node[flowlbl, pos=0.5, right=2pt]{não};
  \draw[flow] (a4) -- (a5);
  \draw[flow] (a5) -- (f);
\end{tikzpicture}}
\caption[AT-07 Conversa com agente]{\req{AT-07} --- conversa com o agente. A
  ferramenta de busca documental aparece em laranja porque é ela que carrega o
  \textit{guardrail}: o modelo recebe um aviso estruturado, não a liberdade de
  decidir se existe procedimento.}
\label{fig:at07}
\end{figure}

% ---------------------------------------------------------------------------
\section{AT-08 --- Ingestão de telemetria industrial}

\begin{landscape}
\begin{figure}[H]
\centering
\fitpage{%
\begin{tikzpicture}[x=1cm, y=1cm]

  \foreach \y/\h/\nome in {0/2.4/{\textit{Gateway} / simulador},
                           -2.4/5.4/{\textit{Worker} de ingestão},
                           -7.8/2.4/{\textit{Broker} MQTT},
                           -10.2/2.2/{Recursos}} {
    \draw[lane] (3.2,\y) rectangle (23.0,\y-\h);
    \draw[lanehdr] (0.4,\y) rectangle (3.2,\y-\h);
    \node[font=\scriptsize\bfseries\color{white}, text width=24mm]
      at (1.8,\y-\h/2) {\nome};
  }

  \node[actstart] (i) at (4.6,-1.2) {};
  \node[action, text width=32mm] (p1) at (8.0,-1.2)
    {Publicar em \cd{sensors/\{maq\}/telemetry} com QoS 1};

  \node[action, text width=30mm] (w1) at (8.0,-3.6)
    {Extrair identificador da máquina do tópico};
  \node[decision, text width=22mm] (d1) at (12.4,-3.6) {JSON e esquema válidos?};
  \node[actionerr, text width=30mm] (dlq) at (17.6,-3.6)
    {Compor registro de rejeição: motivo + \textit{payload}};
  \node[action, text width=30mm] (w2) at (8.0,-6.0)
    {Canonicalizar a anotação, se houver};
  \node[action, text width=30mm] (w3) at (12.4,-6.0)
    {\cd{merge} idempotente em \cd{events}};
  \node[action, text width=30mm] (w4) at (16.8,-6.0)
    {Chamar \cd{POST /diagnose}};
  \node[decision, text width=22mm] (d2) at (21.0,-6.0) {É falha?};

  \node[action, text width=30mm] (b1) at (17.6,-9.0)
    {Publicar em \cd{.../alerts} (QoS 1)};
  \node[action, text width=30mm] (b2) at (11.0,-9.0)
    {Publicar em \cd{.../dlq} (QoS 1)};
  \node[actend] (f) at (5.4,-9.0) {};

  \node[actionsys, text width=30mm] (r1) at (12.4,-11.3)
    {PostgreSQL --- histórico};
  \node[actionsys, text width=30mm] (r2) at (17.6,-11.3)
    {SCADA / CMMS --- assinante};

  \draw[flow] (i) -- (p1);
  \draw[flow] (p1.south) -- (8.0,-2.9) -- (w1.north);
  \draw[flow] (w1) -- (d1);
  \draw[flow] (d1) -- (dlq) node[flowlbl, midway, above=1pt]{não};
  \draw[flow] (dlq.south) -- (17.6,-7.2) -- (11.0,-7.2) -- (b2.north);
  \draw[flow] (d1.south) -- (12.4,-4.6) -- (8.0,-4.6) -- (w2.north)
    node[flowlbl, pos=0.32, above=1pt]{sim};
  \draw[flow] (w2) -- (w3);
  \draw[flow] (w3) -- (w4);
  \draw[flow] (w4) -- (d2);
  \draw[flow] (d2.south) -- (21.0,-7.2) -- (17.6,-7.2) -- (b1.north)
    node[flowlbl, pos=0.22, above=1pt]{sim};
  \node[actend] (f2) at (22.2,-2.9) {};
  \draw[flow] (d2.north) -- (21.0,-4.6) -- (22.2,-4.6) -- (f2)
    node[flowlbl, pos=0.55, above=1pt]{não: apenas persiste};
  \draw[flow] (b2) -- (f);
  \draw[flow] (w3.south) -- (12.4,-10.4);
  \draw[flow] (b1.south) -- (17.6,-10.4);

\end{tikzpicture}}
\caption[AT-08 Ingestão MQTT]{\req{AT-08} --- ingestão de telemetria industrial.
  A persistência precede o diagnóstico: perder o diagnóstico é recuperável,
  perder o dado bruto não é. O ramo de rejeição nunca descarta em silêncio.}
\label{fig:at08}
\end{figure}
\end{landscape}

% ---------------------------------------------------------------------------
\section{AT-09 --- Cadastro de documento pela interface}

\begin{figure}[H]
\centering
\fitwidth{%
\begin{tikzpicture}[x=1cm, y=1cm]
  \node[actstart] (i) at (3.4,0) {};
  \node[action] (a1) at (3.4,-1.4) {Usuário abre a seção \textbf{Documentos}};
  \node[action] (a2) at (3.4,-3.1) {Interface lista documentos e cobertura
    corrente};
  \node[action] (a3) at (3.4,-4.9) {Usuário anexa o PDF, informa título e
    famílias};
  \node[decision] (d1) at (3.4,-7.0) {Campos preenchidos?};
  \node[actionerr] (e1) at (11.0,-7.0) {Exibir erro no formulário};
  \node[action] (a4) at (3.4,-9.3) {Enviar \cd{POST /documents}
    (\cd{multipart})};
  \node[decision] (d2) at (3.4,-11.4) {Indexação bem-sucedida?};
  \node[actionerr] (e2) at (11.0,-11.4) {Exibir código e detalhe do erro
    {\scriptsize (família inválida, PDF vazio)}};
  \node[actionok] (a5) at (3.4,-13.7) {Exibir contagem de \textit{chunks} e
    famílias atendidas};
  \node[action] (a6) at (3.4,-15.5) {Invalidar o \textit{cache} de consultas da
    interface};
  \node[actionok] (a7) at (3.4,-17.3) {A família passa a ser coberta na próxima
    requisição};
  \node[actend] (f) at (3.4,-19.1) {};

  \draw[flow] (i) -- (a1);
  \draw[flow] (a1) -- (a2);
  \draw[flow] (a2) -- (a3);
  \draw[flow] (a3) -- (d1);
  \draw[flow] (d1) -- (e1) node[flowlbl, midway, above=1pt]{não};
  \draw[flow] (e1.south) -- (11.0,-4.9) -- (a3.east);
  \draw[flow] (d1) -- (a4) node[flowlbl, pos=0.5, right=2pt]{sim};
  \draw[flow] (a4) -- (d2);
  \draw[flow] (d2) -- (e2) node[flowlbl, midway, above=1pt]{não};
  \draw[flow] (d2) -- (a5) node[flowlbl, pos=0.5, right=2pt]{sim};
  \draw[flow] (a5) -- (a6);
  \draw[flow] (a6) -- (a7);
  \draw[flow] (a7) -- (f);
\end{tikzpicture}}
\caption[AT-09 Cadastro de documento]{\req{AT-09} --- cadastro de documento pela
  interface. O ciclo se fecha em uma única sessão de trabalho: o engenheiro
  descobre a lacuna no diagnóstico, escreve o procedimento, cadastra e o sistema
  passa a prescrever.}
\label{fig:at09}
\end{figure}

% ---------------------------------------------------------------------------
\section{AT-10 --- Indexação automática no primeiro \textit{boot}}

\begin{figure}[H]
\centering
\fitwidth{%
\begin{tikzpicture}[x=1cm, y=1cm]
  \node[actstart] (i) at (3.4,0) {};
  \node[decision] (d0) at (3.4,-1.9) {\cd{RAG\_BOOTSTRAP} habilitado?};
  \node[actionsys] (n0) at (11.2,-1.9) {Não fazer nada};
  \node[decision] (d1) at (3.4,-4.4) {Índice já contém famílias?};
  \node[actionsys] (n1) at (11.2,-4.4) {Registrar em \textit{log}:
    \enquote{bootstrap dispensado}};
  \node[decision] (d2) at (3.4,-6.9) {Índice inspecionável?};
  \node[actionerr] (n2) at (11.2,-6.9) {Advertir e cancelar --- índice ausente ou
    corrompido};
  \node[action] (a1) at (3.4,-9.2) {Iniciar \textbf{\textit{thread} separada}
    (\textit{daemon})};
  \node[actionok] (a2) at (3.4,-11.0) {\textit{Startup} conclui --- o
    \textit{healthcheck} passa imediatamente};
  \node[action] (a3) at (11.2,-11.0) {Indexar cada PDF declarado na cobertura
    inicial (AT-04)};
  \node[action] (a4) at (11.2,-13.0) {Registrar contagem de \textit{chunks} por
    documento};
  \node[actend] (f) at (3.4,-13.0) {};
  \node[actend] (f2) at (11.2,-14.8) {};

  \draw[flow] (i) -- (d0);
  \draw[flow] (d0) -- (n0) node[flowlbl, midway, above=1pt]{não};
  \draw[flow] (d0) -- (d1) node[flowlbl, pos=0.5, right=2pt]{sim};
  \draw[flow] (d1) -- (n1) node[flowlbl, midway, above=1pt]{sim};
  \draw[flow] (d1) -- (d2) node[flowlbl, pos=0.5, right=2pt]{não};
  \draw[flow] (d2) -- (n2) node[flowlbl, midway, above=1pt]{não};
  \draw[flow] (d2) -- (a1) node[flowlbl, pos=0.5, right=2pt]{sim};
  \draw[flow] (a1) -- (a2);
  \draw[flow] (a1.east) -- (11.2,-9.2) -- (a3.north)
    node[flowlbl, pos=0.35, above=1pt]{assíncrono};
  \draw[flow] (a2) -- (f);
  \draw[flow] (a3) -- (a4);
  \draw[flow] (a4) -- (f2);

  \node[actnote, text width=50mm, anchor=north west] at (16.0,-9.4)
    {A \textit{thread} separada é obrigatória: a indexação faz dezenas de
     chamadas ao serviço de linguagem --- incluindo OCR --- e leva minutos.
     Bloquear o \textit{startup} faria o orquestrador reprovar a instância antes
     de a indexação terminar, em laço infinito.};
  \draw[actnoteline] (15.9,-10.8) -- (13.4,-10.8);
\end{tikzpicture}}
\caption[AT-10 Bootstrap documental]{\req{AT-10} --- indexação automática no
  primeiro \textit{boot}. Três guardas antecedem a execução, e a indexação
  ocorre fora do caminho de \textit{startup}.}
\label{fig:at10}
\end{figure}

% ---------------------------------------------------------------------------
\section{AT-11 --- Consulta a um painel analítico}

\begin{figure}[H]
\centering
\fitwidth{%
\begin{tikzpicture}[x=1cm, y=1cm]
  \node[actstart] (i) at (3.4,0) {};
  \node[action] (a1) at (3.4,-1.4) {Usuário escolhe seção e subpainel};
  \node[decision] (d0) at (3.4,-3.4) {API acessível?};
  \node[actionerr] (e0) at (11.0,-3.4) {Interromper com o endereço configurado
    na mensagem};
  \node[actend] (fe) at (11.0,-5.4) {};
  \node[action] (a2) at (3.4,-5.8) {Requisitar \textbf{apenas} os dados do
    painel ativo};
  \node[decision] (d1) at (3.4,-7.9) {Agregação disponível?};
  \node[actionerr] (e1) at (11.0,-7.9) {Exibir aviso no painel e manter os
    demais gráficos};
  \node[action] (a3) at (3.4,-10.2) {Aplicar filtros locais (rotação, classe de
    máquina, modo)};
  \node[action] (a4) at (3.4,-12.1) {Aplicar a paleta validada e o
    \textit{chrome} comum};
  \node[actionok] (a5) at (3.4,-14.0) {Renderizar gráficos com legenda
    explicativa};
  \node[actend] (f) at (3.4,-15.8) {};

  \draw[flow] (i) -- (a1);
  \draw[flow] (a1) -- (d0);
  \draw[flow] (d0) -- (e0) node[flowlbl, midway, above=1pt]{não};
  \draw[flow] (e0) -- (fe);
  \draw[flow] (d0) -- (a2) node[flowlbl, pos=0.5, right=2pt]{sim};
  \draw[flow] (a2) -- (d1);
  \draw[flow] (d1) -- (e1) node[flowlbl, midway, above=1pt]{não};
  \draw[flow] (e1.south) -- (11.0,-10.2) -- (a3.east);
  \draw[flow] (d1) -- (a3) node[flowlbl, pos=0.5, right=2pt]{sim};
  \draw[flow] (a3) -- (a4);
  \draw[flow] (a4) -- (a5);
  \draw[flow] (a5) -- (f);
\end{tikzpicture}}
\caption[AT-11 Consulta a painel]{\req{AT-11} --- consulta a um painel
  analítico. A indisponibilidade de uma agregação degrada apenas o gráfico
  afetado; o painel continua útil.}
\label{fig:at11}
\end{figure}

% ---------------------------------------------------------------------------
\section{AT-12 --- Integração contínua}

\begin{figure}[H]
\centering
\fitwidth{%
\begin{tikzpicture}[x=1cm, y=1cm]
  \node[actstart] (i) at (1.6,0) {};
  \node[action, text width=28mm] (a1) at (4.6,0)
    {\textit{Push} em \cd{main} ou \textit{pull request}};
  \node[action, text width=28mm] (a2) at (8.8,0)
    {\textbf{\textit{Lint}}\newline{\scriptsize \cd{ruff check} +
     \cd{ruff format --check}}};
  \node[decision, text width=20mm] (d1) at (12.8,0) {Passou?};
  \node[action, text width=30mm] (a3) at (8.8,-3.2)
    {\textbf{Testes}\newline{\scriptsize 63 testes; Postgres como serviço;
     linguagem substituída por \textit{fake}}};
  \node[decision, text width=20mm] (d2) at (13.4,-3.2) {Passou?};
  \node[action, text width=28mm] (a4) at (8.8,-6.4)
    {\textbf{\textit{Build}}\newline{\scriptsize \cd{docker build} da imagem
     única}};
  \node[decision, text width=20mm] (d3) at (12.8,-6.4) {Passou?};
  \node[actionok, text width=30mm] (a5) at (8.8,-9.4)
    {Publicação autorizada --- a nuvem só implanta após o CI aprovar};
  \node[actionerr, text width=26mm] (e) at (17.4,-3.2)
    {Interromper e reportar no \textit{pull request}};
  \node[actend] (fe) at (17.4,-5.4) {};
  \node[actend] (f) at (8.8,-11.2) {};

  \draw[flow] (i) -- (a1);
  \draw[flow] (a1) -- (a2);
  \draw[flow] (a2) -- (d1);
  \draw[flow] (d1.south) -- (12.8,-1.7) -- (8.8,-1.7) -- (a3.north)
    node[flowlbl, pos=0.32, above=1pt]{sim};
  \draw[flow] (a3) -- (d2);
  \draw[flow] (d2.south) -- (13.4,-4.9) -- (8.8,-4.9) -- (a4.north)
    node[flowlbl, pos=0.32, above=1pt]{sim};
  \draw[flow] (a4) -- (d3);
  \draw[flow] (d3.south) -- (12.8,-8.0) -- (8.8,-8.0) -- (a5.north)
    node[flowlbl, pos=0.32, above=1pt]{sim};
  \draw[flow] (a5) -- (f);
  \draw[flow] (d1.east) -- (15.4,0) -- (15.4,-2.6) -- (e.north)
    node[flowlbl, pos=0.28, above=1pt]{não};
  \draw[flow] (d2.east) -- (e.west) node[flowlbl, midway, above=1pt]{não};
  \draw[flow] (d3.east) -- (15.4,-6.4) -- (15.4,-3.8) -- (e.south)
    node[flowlbl, pos=0.28, below=1pt]{não};
  \draw[flow] (e) -- (fe);
\end{tikzpicture}}
\caption[AT-12 Integração contínua]{\req{AT-12} --- integração contínua. As
  etapas são sequenciais e bloqueantes: um apontamento de estilo impede a
  execução dos testes, que por sua vez precedem a construção da imagem.
  Nenhum teste do fluxo chama o serviço de linguagem real.}
\label{fig:at12}
\end{figure}

% ---------------------------------------------------------------------------
\section{AT-13 --- Provisionamento e implantação em nuvem}

\begin{figure}[H]
\centering
\fitwidth{%
\begin{tikzpicture}[x=1cm, y=1cm]
  \node[actstart] (i) at (3.4,0) {};
  \node[action] (a1) at (3.4,-1.4) {Instalar o SDK de infraestrutura e
    autenticar};
  \node[action] (a2) at (3.4,-3.1) {Inicializar o projeto};
  \node[action] (a3) at (3.4,-4.8) {\textbf{Planejar}: revisar o \textit{diff}
    da topologia};
  \node[decision] (d1) at (3.4,-6.8) {\textit{Diff} aceitável?};
  \node[actionerr] (e1) at (11.0,-6.8) {Ajustar a declaração e replanejar};
  \node[action] (a4) at (3.4,-9.1) {\textbf{Aplicar}: criar banco, serviços,
    volume e grupos};
  \node[action] (a5) at (3.4,-10.9) {Definir segredos por entrada padrão
    {\scriptsize (fora do histórico do \textit{shell})}};
  \node[action] (a6) at (3.4,-12.8) {Expor domínios HTTPS e o \textit{proxy} TCP
    do \textit{broker}};
  \node[action] (a7) at (3.4,-14.7) {Carregar o histórico da máquina local
    apontando ao banco remoto};
  \node[action] (a8) at (3.4,-16.6) {Aguardar a indexação automática
    (AT-10)};
  \node[decision] (d2) at (3.4,-18.7) {Verificação pós-implantação passou?};
  \node[actionerr] (e2) at (11.0,-18.7) {Inspecionar \textit{logs} e variáveis
    efetivas};
  \node[actionok] (a9) at (3.4,-21.0) {Sistema operacional em nuvem};
  \node[actend] (f) at (3.4,-22.8) {};

  \draw[flow] (i) -- (a1);
  \draw[flow] (a1) -- (a2);
  \draw[flow] (a2) -- (a3);
  \draw[flow] (a3) -- (d1);
  \draw[flow] (d1) -- (e1) node[flowlbl, midway, above=1pt]{não};
  \draw[flow] (e1.north) -- (11.0,-4.8) -- (a3.east);
  \draw[flow] (d1) -- (a4) node[flowlbl, pos=0.5, right=2pt]{sim};
  \draw[flow] (a4) -- (a5);
  \draw[flow] (a5) -- (a6);
  \draw[flow] (a6) -- (a7);
  \draw[flow] (a7) -- (a8);
  \draw[flow] (a8) -- (d2);
  \draw[flow] (d2) -- (e2) node[flowlbl, midway, above=1pt]{não};
  \draw[flow] (e2.south) -- (11.0,-21.0) -- (a9.east);
  \draw[flow] (d2) -- (a9) node[flowlbl, pos=0.5, right=2pt]{sim};
  \draw[flow] (a9) -- (f);
\end{tikzpicture}}
\caption[AT-13 Provisionamento em nuvem]{\req{AT-13} --- provisionamento e
  implantação. O passo de planejamento é obrigatório e revisável: a topologia é
  declarada em código, e nenhuma mudança é aplicada sem que o \textit{diff}
  tenha sido lido.}
\label{fig:at13}
\end{figure}

% ---------------------------------------------------------------------------
\section{AT-14 --- Ciclo de tratamento de lacuna documental}

Este é o único diagrama de nível de negócio: ele descreve o ciclo completo que
o sistema habilita, atravessando pessoas e \textit{software}. É a razão de ser
da solução.

\begin{landscape}
\begin{figure}[H]
\centering
\fitpage{%
\begin{tikzpicture}[x=1cm, y=1cm]

  \foreach \y/\h/\nome in {0/2.6/{Planta / sensor},
                           -2.6/2.8/{Sistema},
                           -5.4/3.0/{Técnico de manutenção},
                           -8.4/3.2/{Engenheiro de manutenção}} {
    \draw[lane] (3.2,\y) rectangle (23.0,\y-\h);
    \draw[lanehdr] (0.4,\y) rectangle (3.2,\y-\h);
    \node[font=\scriptsize\bfseries\color{white}, text width=24mm]
      at (1.8,\y-\h/2) {\nome};
  }

  \node[actstart] (i) at (4.6,-1.3) {};
  \node[action, text width=30mm] (p1) at (8.2,-1.3)
    {Vibração anômala detectada pelo sensor};

  \node[action, text width=30mm] (s1) at (8.2,-4.0)
    {Diagnóstico: família prevista + ocorrências semelhantes};
  \node[decision, text width=22mm] (d1) at (12.8,-4.0) {Documentada?};
  \node[actionok, text width=30mm] (s2) at (17.4,-4.0)
    {Prescrição com citação de documento, seção e página};
  \node[actionerr, text width=30mm] (s3) at (21.4,-2.9)
    {Aviso + sugestão de registro};

  \node[action, text width=30mm] (t1) at (17.4,-6.8)
    {Executar o procedimento prescrito};
  \node[actionok, text width=30mm] (t2) at (12.4,-6.8)
    {Registrar conclusão e validar critérios de aceitação};
  \node[actend] (f1) at (8.0,-6.8) {};

  \node[action, text width=32mm] (e1) at (21.0,-9.8)
    {Priorizar a lacuna no painel {\scriptsize (frequência $\times$ severidade)}};
  \node[action, text width=32mm] (e2) at (15.6,-9.8)
    {Redigir o procedimento técnico};
  \node[action, text width=32mm] (e3) at (10.4,-9.8)
    {Cadastrar o PDF (AT-09)};
  \node[actionok, text width=30mm] (e4) at (5.6,-9.8)
    {Família passa a ser coberta};

  \draw[flow] (i) -- (p1);
  \draw[flow] (p1.south) -- (8.2,-3.1) -- (s1.north);
  \draw[flow] (s1) -- (d1);
  \draw[flow] (d1) -- (s2) node[flowlbl, midway, above=1pt]{sim};
  \draw[flow] (d1.north) -- (12.8,-2.9) -- (s3.west)
    node[flowlbl, pos=0.4, above=1pt]{não};
  \draw[flow] (s2.south) -- (17.4,-5.5) -- (t1.north);
  \draw[flow] (t1) -- (t2);
  \draw[flow] (t2) -- (f1);
  \draw[flow] (s3.south) -- (21.4,-8.5) -- (21.0,-8.5) -- (e1.north);
  \draw[flow] (e1) -- (e2);
  \draw[flow] (e2) -- (e3);
  \draw[flow] (e3) -- (e4);
  \draw[flow] (e4.west) -- (3.9,-9.8) -- (3.9,-4.0) -- (s1.west)
    node[flowlbl, pos=0.62, right=2pt]{o próximo diagnóstico já prescreve};

\end{tikzpicture}}
\caption[AT-14 Ciclo de lacuna documental]{\req{AT-14} --- ciclo de tratamento
  de lacuna documental. O laço de retorno é o valor central da solução: uma
  falha sem procedimento deixa de ser um beco sem saída e passa a ser uma tarefa
  priorizada, que ao ser concluída realimenta o sistema.}
\label{fig:at14}
\end{figure}
\end{landscape}
